\question
Неизвестный алхимик средневековья решил придумать новое зелье для короля, но столкнулся с проблемой, что не все ингредиенты можно смешивать, а количество котлов и колб ограничено. Изучив древние письмена, он составил список ингредиентов, которые нельзя варить и хранить вместе:

\begin{itemize}
\item Глаз змеи – перо феникса;

\item Родниковая вода – корень мандрагоры;

\item Первоцвет – кристалл воздуха;

\item Первоцвет – глициния;

\item Глаз змеи - родниковая вода;

\item Зуб дракона - первоцвет;

\item Белладонна – корень мандрагоры;

\item Белладонна - глициния;

\item Глициния – корень мандрагоры;

\item Зуб дракона - кристалл воздуха;
\end{itemize}

Используя теорию графов, помогите алхимику вычислить количество котлов и колб, требующееся для изготовления эликсиров из всех вышеуказанных ингредиентов. 

---------------

Автор -- Елизавета Котельникова